\chapter{Discussion}\label{chap:discussion}
\section{Introduction}
The primary goal of this chapter is to contextualize the empirical findings of this study within the theoretical framework presented in Chapter \ref{chap:SocialScience}. Grounded in Eccles' Expectancy-Value Theory (\cite{eccles1989}) and Bandura's Social Cognitive Theory (\cite{Bandura2001}), this research investigates whether a narrative-driven seminar focusing on female role models can actively challenge the cultural stereotype of mathematics as a male domain (\cite{Cheryan2015,Nosek2002}). To achieve this, a quantitative within-group design was employed, tracking changes across four core psychological constructs applied to mathematics: stereotypes, self-efficacy, anxiety, and attitudes.

Rather than analyzing individual variables in isolation, this discussion aims to unpack the underlying psychological mechanisms that explain how and why early adolescents react to exposure to female role models. By evaluating each research question sequentially, the following sections link the statistical outcomes to the socio-cognitive literature. Additionally, this chapter outlines concrete pedagogical and didactical recommendations, and concludes with an analysis of the methodological limitations that restrict the study's generalizability, paving the way for future longitudinal research.

\section{Research Questions}
\subsection{Gender $\times$ Stereotype (RQ1)}
The first research question explored the degree to which the school-based seminar could diminish explicit gender stereotypes regarding mathematics among lower secondary school students. At the baseline ($T_0$), the scores for explicit stereotypes were notably low and statistically comparable between boys and girls ($d = 0.01$). This absence of an initial gender gap suggests that, in this particular educational setting, early adolescents do not openly support the idea of male superiority in mathematics. This finding is consistent with recent European trends that indicate a decrease in overt sexism within school environments (\cite{Master2020}).

However, the qualitative data gathered during the interactive introduction reveals a more complex perspective on implicit or unexamined biases. When asked to anonymously list famous mathematicians, students predominantly named only men, while the few women mentioned were largely associated with broader scientific fields—like Marie Curie and Mileva Marić who was primarily recognized through her relationship and collaboration with Albert Einstein. Additionally, when students noticed that their lists included only men, they initially justified this gap by suggesting that women in the past ``weren't interested'' in math or ``didn't like it.''

The contrast between low explicit stereotype scores and a predominantly male cognitive mapping highlights what Greenwald et al. (\citeyear{Greenwald2002}) refer to as a structural absence of counter-stereotypical examples in memory. Gunderson et al. (\citeyear{Gunderson2012}) argue that children absorb cultural norms associated with mathematical fields without necessarily agreeing with overt expressions of discrimination. The seminar's focus on the historical journeys of figures like Hypatia, Maria Gaetana Agnesi, Sophie Germain, and Sofya Kovalevskaya actively challenged the default ``male genius'' narrative (\cite{Leslie2015}). By clearly addressing the institutional and legal obstacles that compelled women to use male pseudonyms or study in secrecy, the seminar reshaped the students' understanding: the historical scarcity of women was reframed not as a natural lack of interest, but as a result of systemic exclusion (\cite{Cheryan2015}).

This cognitive re-calibration explains the highly significant main effect of Time, which revealed a substantial reduction in stereotype endorsement from pre-test to post-test across the entire sample ($F(1, 349) = 33.63, p < .001, \eta_p^2 = 0.088$), followed by a notable effect size ($d = -0.31$).
Additionally, the absence of a Time $\times$ Gender interaction ($F(1, 349) = 0.135, p = .713$) highlights that boys showed were equally receptive to this narrative shift.
Rather than provoking a defensive or distancing response, introducing male students to figures such as Maryam Mirzakhani broadened their understanding of mathematical identity. This consistent change supports the argument made by Smith and Farkas (\citeyear{Smith2023}) that confronting the masculine portrayal of math within peer groups reduces the common reinforcement of traditional roles, fostering a more equitable classroom environment for both genders equally.

\subsection{Gender $\times$ Self-Efficacy (RQ2)}
The second research question evaluated how the seminar influenced students' academic self-efficacy, which is defined as their personal confidence in performing specific mathematical tasks successfully. At the initial measurement ($T_0$), a significant gender gap was observed: even though female students obtained equivalent or better grades in standard lower secondary school tracks, they reported significantly lower self-efficacy compared to their male counterparts ($t(349) = -2.25, p = .025, d = -0.24$). This initial difference aligns closely with the traditional socio-cognitive patterns identified by Hackett and Betz (\citeyear{HACKETT1981326}), where gender-role socialization leads girls to underestimate their mathematical skills, despite their academic performance being on par with that of boys (\cite{Pajares2005}).

After the intervention ($T_1$), the mixed ANOVA indicated that the main effect of Time on mathematics self-efficacy did not reach statistical significance ($F(1, 340)= 3.11, p= .079$). However, a marginally positive trend was observed in the paired t-test ($t(341)=1.78, p= .076, d= 0.10$). This resistance to short-term change can be effectively explained by Bandura's (\citeyear{Bandura1997}) hierarchical model of self-efficacy sources. According to Social Cognitive Theory, self-efficacy primarily arises from Authentic Mastery Experiences—the direct, personal feedback loop of success and failure in completing a task. Vicarious experiences (such as observing role models) and social persuasion (such as encouragement from educators or peers) play secondary, less stable roles (\cite{Usher2008}).

Although the seminar offered strong vicarious exposure through the presentation of resilient female role models, a 60-minute narrative intervention cannot deliver immediate personal mastery experiences, as the students did not engage in solving complex mathematical problems or receive evaluative feedback. As a result, while the girls' perceptions of who is capable of doing math broadened—evident from the significant reduction in explicit stereotypes—their self-assessment of their own mathematical abilities remained stable.

The gap between changing stereotypes and improving self-efficacy corresponds with the longitudinal research conducted by Zeldin et al. (\citeyear{Zel}), which demonstrates that fostering strong self-efficacy in women necessitates ongoing relational support and continuous opportunities for personal achievement in an inclusive classroom setting. However, the slightly positive trend suggests that vicarious experiences could act as an important psychological trigger, preparing students to confront future mathematical challenges with greater confidence; yet, this must be accompanied by focused mastery experiences to create lasting structural changes.

\subsection{Gender $\times$ Mathematics Anxiety (RQ3)}
Improve it
The third research question examined the intervention's ability to reduce levels of math anxiety. At the baseline measurement ($T_0$), there was a significant gender disparity in math anxiety ($t(339) = 4.91, p < .001, d = 0.53$), with female students reporting considerably higher levels of anxiety related to learning and testing situations. This notable gap aligns with the theoretical mechanisms of Stereotype Threat discussed by Spencer et al. (\citeyear{Spencer1999}). Even without overt or conscious endorsement of sexism, girls often find themselves in an intellectual environment filled with implicit performance anxieties. This psychological pressure results in heightened vigilance and self-monitoring, which can disrupt working memory and increase somatic anxiety (\cite{Beilock2004, Ramirez2013}).

After the intervention, a significant main effect of Time was observed, indicating a consistent decrease in math anxiety across the entire sample ($F(1, 330) = 12.55, p < .001, \eta_p^2 = 0.037, d = -0.19$). Importantly, the interaction between Time and Gender was found to be non-significant ($F(1, 330) = 0.59, p = .445$), suggesting that the seminar was equally effective for both male and female students.

To understand this phenomenon, we need to examine the specific narrative choices made during the presentation. Instead of portraying these mathematicians as infallible or biologically predetermined geniuses, the focus was on their human vulnerabilities and academic challenges. For instance, students were particularly engaged by the biographical detail that Maryam Mirzakhani was initially discouraged by a teacher who told her she lacked mathematical aptitude.

By framing mathematical learning as a process of resilience and strategy instead of a fixed, innate ability, the seminar transformed the meaning of difficulty. According to Schmader et al. \citeyear{Schmader2008}), math anxiety increases when individuals view errors or confusion as evidence of an inherent deficiency (such as thinking, ``I'm a girl, so I can't do this'' or ``I'm not smart enough''). By de-stigmatizing struggle and connecting advanced mathematics to practical, creative visual strategies, the seminar removed the rigid and exclusionary image often associated with the discipline. For girls, this approach directly undermined the motivational path of stereotype anxiety by making the subject more relatable (\cite{Nosek2002}). For boys, it likely reduced the pressure tied to the expectations, which can alleviate general test anxiety and improve their emotional connections to the subject.

\subsection{Gender $\times$ Attitudes Toward Mathematics (RQ4)}
The fourth research question evaluated changes in students' overall attitudes toward mathematics, focusing on their interest, enjoyment, and perceived everyday utility. Baseline measures ($T_0$) highlighted a significant gender gap, with boys displaying notably more positive attitudes than girls ($t(347) = -3.12, p = .002, d = -0.33$). This finding is consistent with the longitudinal patterns identified by Eccles (\citeyear{eccles1989}), which show that girls' subjective task value and interest in mathematics begin to decline sharply during early adolescence. According to Expectancy-Value Theory (\cite{EcclesWigfield2020, Eccles1983}), this decline occurs because girls increasingly view advanced mathematics as unrelated to interpersonal relationships, isolating, and not aligned with their developing self-concept.

The post-intervention analysis revealed a statistically significant main effect of Time, indicating an overall improvement in general attitudes across the entire sample ($F(1, 339) = 4.24, p = .040, \eta_p^2 = 0.012, d = 0.11$). Importantly, the absence of a significant Time $\times$ Gender interaction ($p = .300$) confirmed that both genders experienced a similar increase in a positive direction.

This positive shift in attitude can be directly linked to the contextualization strategies used in the seminar. Instead of presenting advanced mathematical topics as dry, abstract concepts or mechanical operations—common perceptions among lower secondary school students—the researcher explicitly connected Mirzakhani's work on the geometry of curved surfaces to tangible, modern technologies like Global Positioning Systems (GPS) and computer graphics. According to the Expectancy-Value model, demonstrating the practical applications of a discipline increases its perceived utility, which is a crucial factor in fostering intrinsic motivation (\cite{Guiso2008}).

Furthermore the biographical narrative of Maria Gaetana Agnesi, a polyglot who mastered multiple languages and differential calculus, challenged the immage of the isolated genius often associated with STEM personalities (\cite{Diekman2010}).
Her story shows that world-class mathematicians can successfully balance their literary interests, family responsibilities, and commitment to humanitarian efforts. This powerful counter-narrative allows both girls and boys to view mathematics not as a restrictive subject but as a creative, expansive, and universally applicable human endeavor, thereby reducing the ``identity cost'' of entering traditionally masculine fields.

\section{Pedagogical Implications and Reccommendations}

Finding that a brief, narrative seminar can simoultanoeusly reduce explicit stereotypes, lower math anxiety, and enhance positive attitudes provides valuable, practical recommendations for educational methods.
To transform these temporary cognitive and emotional changes into lasting academic results, the following teaching strategies are suggested:

\begin{itemize}
    \item \textbf{Integrate ``Herstories'' into the Core Curriculum:}Integrate ``Herstories'' into the Core Curriculum: Instead of limiting the contributions of female scientists to optional modules or designating a single day for gender equity, biographical profiles of women should be incorporated directly into daily mathematics lessons. For instance, when teaching coordinate geometry, educators could introduce Agnesi's curve; when discussing prime numbers, they might highlight Sophie Germain's foundational work on Fermat's Last Theorem. This systematic integration normalizes the contributions of female intellect within the primary historical narrative, providing strong counter-stereotypical examples necessary to reshape students' cognitive understanding (\cite{Greenwald2002}).
    \item  \textbf{Challeng the Myth of Natural Genius:} Educators should actively work against the fixed mindset that views success in mathematics as a product of effortless, inherent ability (\cite{Dweck2008}). By emphasizing that renowned mathematicians faced challenges, sought guidance from mentors, employed innovative strategies, encountered rejection, and consistently reevaluated their objectives, we can create classrooms that are focused on growth. When students recognize that mistakes are essential tools for learning rather than indicators of fundamental inadequacy, the psychological stress associated with stereotype threat and performance anxiety is greatly reduced (\cite{Schmader2008}).
    \item \textbf{Incorporate Visual and Low-Floor, High-Ceiling Activities:} Incorporate Visual and Low-Floor, High-Ceiling Activities: Emulating Maryam Mirzakhani's intuitive approach, teachers should foster visual intuition, geometric sketch-mapping, and creative problem-solving in the classroom. Utilizing cooperative, hands-on strategies and ``low-floor, high-ceiling'' activities promotes an inclusive learning environment. These methods minimize intense public competition and reduce the reliance on time-pressured drills, which can increase math anxiety, particularly among female students, while also alienating struggling learners of all genders (\cite{Boaler2016}).
    \item \textbf{Career Counseling:}Active Career and Context Counseling: Educators should clearly illustrate the connections between lower secondary school mathematics and exciting, real-world career paths. By demonstrating how advanced geometric and algebraic concepts support modern technologies such as artificial intelligence, robotics, GPS algorithms, and climate modeling, we can enhance the relevance of mathematics under the Expectancy-Value model (\cite{EcclesWigfield2020}). This contextualization aids students—particularly girls—in overcoming the ``identity cost'' associated with STEM fields. It encourages them to make informed, non-stereotypical choices regarding their academic tracks, ultimately leading to the selection of advanced courses later in their educational journeys.
\end{itemize}

\section{Limitations and Future Research Directions}
While this study provides important preliminary evidence for the effectiveness of the intervention, it is essential to recognize several methodological limitations in order to accurately interpret the results. The primary structural limitation lies in the use of a quantitative within-group quasi-experimental design. Without a randomized control group participating in an alternate or neutral activity, we cannot fully determine the specific effects of the intervention, as it leaves room for influences such as testing effects, maturation, or social desirability biases. For instance, students might have reported lower levels of stereotype endorsement or less anxiety at $T_1$ simply due to their awareness of the researcher's expectations. Therefore, future studies should adopt a rigorous randomized controlled trial (RCT) design that includes an active control group—like a seminar on literature or general history—to eliminate these confounding factors and more clearly establish causal relationships.

Additionally, the geographical scope and sample limitations affect the immediate generalizability of the findings. The sample was taken exclusively from a single middle school located in the Province of Trieste. Although the sample size was statistically substantial ($N = 362$), the specific socio-demographic and cultural characteristics of this urban student population limit the external validity of the results. To establish broader applicability, it is crucial to test and replicate the intervention in diverse rural, suburban, and varying socio-economic school districts. This step is essential to assess its scalability, robustness, and cultural relevance.

Another significant limitation of the study is its short-term evaluative timeframe, as the post-test questionnaire was administered precisely one week after the seminar. While this approach effectively captured immediate cognitive adjustments and temporary emotional relief, it does not assess the longevity of these psychological changes. Future research should adopt a comprehensive longitudinal design, monitoring students at longer intervals, for example after three months, six months, and one year. This long-term tracking is essential to determine if these emotional changes lead to lasting behavioral outcomes and tangible academic decisions, such as enrolling in scientific high school programs (liceo scientifico).

Finally, this study was based solely on self-reported Likert-scale questionnaires, which mainly assess conscious beliefs. The significant gap observed between low explicit stereotype scores and the predominantly male qualitative listing task highlights that implicit biases function through different, less accessible cognitive channels. To address this measurement limitation, future research should include implicit measures, like the Implicit Association Test (IAT), in combination with longitudinal qualitative interviews. Utilizing both methods would provide a more thorough and nuanced understanding of identity development, the flexibility of self-concept, and the evolution of relational support over time.